
\documentclass{article}
\usepackage{colortbl}
\usepackage{makecell}
\usepackage{multirow}
\usepackage{supertabular}

\begin{document}

\newcounter{utterance}

\twocolumn

{ \footnotesize  \setcounter{utterance}{1}
\setlength{\tabcolsep}{0pt}
\begin{supertabular}{c@{$\;$}|p{.15\linewidth}@{}p{.15\linewidth}p{.15\linewidth}p{.15\linewidth}p{.15\linewidth}p{.15\linewidth}}

    \# & $\;$A & \multicolumn{4}{c}{Game Master} & $\;\:$B\\
    \hline 

    \theutterance \stepcounter{utterance}  

    & & \multicolumn{4}{p{0.6\linewidth}}{\cellcolor[rgb]{0.9,0.9,0.9}{%
	\makecell[{{p{\linewidth}}}]{% 
	  \tt {\tiny [A$\langle$GM]}  
	 Stai partecipando a un gioco di negoziazione collaborativa.\\ \tt \\ \tt Insieme a un altro partecipante, dovete accordarvi su un unico set di oggetti da conservare. Ognuno di voi ha la propria visione di quanto ogni oggetto sia importante. Non sai come l'altro partecipante valuta gli oggetti. Inoltre, ti viene fornito lo sforzo che ogni oggetto richiede.\\ \tt Potete accordarvi su un set solo se lo sforzo totale degli oggetti selezionati non supera un limite condiviso:\\ \tt \\ \tt LIMITE = 3629\\ \tt \\ \tt Ecco i valori di sforzo dei singoli oggetti:\\ \tt \\ \tt Sforzo oggetto = {"A21": 5, "B93": 793, "A11": 21, "B33": 842, "C08": 313, "B39": 842, "B96": 631, "C85": 226, "A82": 87, "B31": 761, "B20": 230, "A60": 287, "B09": 698, "C81": 641, "B21": 882}\\ \tt \\ \tt Ecco la tua visione personale sull'importanza di ogni oggetto:\\ \tt \\ \tt Valori di importanza dell'oggetto = {"A21": 138, "B93": 583, "A11": 868, "B33": 822, "C08": 783, "B39": 65, "B96": 262, "C85": 121, "A82": 508, "B31": 780, "B20": 461, "A60": 484, "B09": 668, "C81": 389, "B21": 808}\\ \tt \\ \tt Obiettivo:\\ \tt \\ \tt Il tuo obiettivo è negoziare un set condiviso di oggetti che ti avvantaggi il più possibile (cioè, massimizzi l'importanza totale per TE), rimanendo entro il LIMITEE. Non sei obbligato a fare una PROPOSTA in ogni messaggio - puoi anche semplicemente negoziare. Tutte le tattiche sono permesse!\\ \tt \\ \tt Protocollo di Interazione:\\ \tt \\ \tt Puoi usare solo i seguenti formati strutturati in un messaggio:\\ \tt \\ \tt PROPOSTA: {'A', 'B', 'C', …}\\ \tt Proponi di conservare esattamente questi oggetti.\\ \tt RIFIUTO: {'A', 'B', 'C', …}\\ \tt Rifiuta esplicitamente la proposta dell'avversario.\\ \tt ARGOMENTAZIONE: {'...'}\\ \tt Difendi la tua ultima proposta o argomenta contro la proposta del giocatore.\\ \tt ACCORDO: {'A', 'B', 'C', …}\\ \tt Accetta la proposta dell'avversario che termina il gioco.\\ \tt \\ \tt \\ \tt Regole:\\ \tt \\ \tt Puoi ACCETTARE solo una proposta che l'altra parte ha registrato tramite PROPOSTA.\\ \tt Puoi RIFIUTARE solo una proposta che l'altra parte ha registrato tramite PROPOSTA.\\ \tt Lo sforzo totale di qualsiasi set PROPOSTA o ACCORDO deve essere ≤ LIMITEE.\\ \tt NON rivelare i tuoi punteggi di importanza nascosti.\\ \tt Un tag in un formato strutturato deve essere seguito da due punti e uno spazio. L'argomento deve essere un set python contenente 0 o più stringhe.\\ \tt Quindi, deve essere della forma TAG: {...}\\ \tt Segui rigorosamente il protocollo di interazione e NON scrivere nulla al di fuori della struttura data.\\ \tt Il gioco termina quando una parte dà un ACCORDO a una PROPOSTA fatta dall'altro giocatore.\\ \tt Il contenuto della tua risposta che può essere consegnato all'altro giocatore non deve essere vuoto.\\ \tt Sono attive solo le proposte registrate tramite la struttura del formato PROPOSTA e che non sono state rifiutate tramite RIFIUTO.\\ \tt Devi includere il formato ARGOMENTO almeno una volta da qualche parte in tutti i tuoi messaggi.\\ \tt  
	  } 
	   } 
	   } 
	 & \\ 
 

    \theutterance \stepcounter{utterance}  

    & \multicolumn{4}{p{0.6\linewidth}}{\cellcolor[rgb]{0.8,1,0.9}{%
	\makecell[{{p{\linewidth}}}]{% 
	  \tt {\tiny [A$\rangle$GM]}  
	 PROPOSTA: {'A11', 'B33', 'C08', 'B31', 'B09', 'A21', 'A82'}\\ \tt \\ \tt ARGOMENTAZIONE: {'Ho selezionato un set ad alto rapporto valore/sforzo che rimane ben entro il limite. Lo sforzo totale è 5 + 21 + 842 + 313 + 761 + 698 + 87 = 2727 ≤ 3629, lasciando margine per eventuali aggiunte mirate. Gli oggetti proposti includono elementi di forte impatto e costo moderato; se hai preferenze su sostituzioni, possiamo valutare uno scambio mantenendo il totale sotto il limite.'} 
	  } 
	   } 
	   } 
	 & & \\ 
 

    \theutterance \stepcounter{utterance}  

    & & \multicolumn{4}{p{0.6\linewidth}}{\cellcolor[rgb]{0.9,0.9,0.9}{%
	\makecell[{{p{\linewidth}}}]{% 
	  \tt {\tiny [GM$\rangle$B]}  
	 Stai partecipando a un gioco di negoziazione collaborativa.\\ \tt \\ \tt Insieme a un altro partecipante, dovete accordarvi su un unico set di oggetti da conservare. Ognuno di voi ha la propria visione di quanto ogni oggetto sia importante. Non sai come l'altro partecipante valuta gli oggetti. Inoltre, ti viene fornito lo sforzo che ogni oggetto richiede.\\ \tt Potete accordarvi su un set solo se lo sforzo totale degli oggetti selezionati non supera un limite condiviso:\\ \tt \\ \tt LIMITE = 3629\\ \tt \\ \tt Ecco i valori di sforzo dei singoli oggetti:\\ \tt \\ \tt Sforzo oggetto = {"A21": 5, "B93": 793, "A11": 21, "B33": 842, "C08": 313, "B39": 842, "B96": 631, "C85": 226, "A82": 87, "B31": 761, "B20": 230, "A60": 287, "B09": 698, "C81": 641, "B21": 882}\\ \tt \\ \tt Ecco la tua visione personale sull'importanza di ogni oggetto:\\ \tt \\ \tt Valori di importanza dell'oggetto = {"A11": 65, "B33": 121, "B21": 138, "C08": 262, "B31": 389, "B09": 461, "B93": 484, "A82": 508, "A60": 583, "B20": 668, "C81": 780, "B96": 783, "A21": 808, "C85": 822, "B39": 868}\\ \tt \\ \tt Obiettivo:\\ \tt \\ \tt Il tuo obiettivo è negoziare un set condiviso di oggetti che ti avvantaggi il più possibile (cioè, massimizzi l'importanza totale per TE), rimanendo entro il LIMITEE. Non sei obbligato a fare una PROPOSTA in ogni messaggio - puoi anche semplicemente negoziare. Tutte le tattiche sono permesse!\\ \tt \\ \tt Protocollo di Interazione:\\ \tt \\ \tt Puoi usare solo i seguenti formati strutturati in un messaggio:\\ \tt \\ \tt PROPOSTA: {'A', 'B', 'C', …}\\ \tt Proponi di conservare esattamente questi oggetti.\\ \tt RIFIUTO: {'A', 'B', 'C', …}\\ \tt Rifiuta esplicitamente la proposta dell'avversario.\\ \tt ARGOMENTAZIONE: {'...'}\\ \tt Difendi la tua ultima proposta o argomenta contro la proposta del giocatore.\\ \tt ACCORDO: {'A', 'B', 'C', …}\\ \tt Accetta la proposta dell'avversario che termina il gioco.\\ \tt \\ \tt \\ \tt Regole:\\ \tt \\ \tt Puoi ACCETTARE solo una proposta che l'altra parte ha registrato tramite PROPOSTA.\\ \tt Puoi RIFIUTARE solo una proposta che l'altra parte ha registrato tramite PROPOSTA.\\ \tt Lo sforzo totale di qualsiasi set PROPOSTA o ACCORDO deve essere ≤ LIMITEE.\\ \tt NON rivelare i tuoi punteggi di importanza nascosti.\\ \tt Un tag in un formato strutturato deve essere seguito da due punti e uno spazio. L'argomento deve essere un set python contenente 0 o più stringhe.\\ \tt Quindi, deve essere della forma TAG: {...}\\ \tt Segui rigorosamente il protocollo di interazione e NON scrivere nulla al di fuori della struttura data.\\ \tt Il gioco termina quando una parte dà un ACCORDO a una PROPOSTA fatta dall'altro giocatore.\\ \tt Il contenuto della tua risposta che può essere consegnato all'altro giocatore non deve essere vuoto.\\ \tt Sono attive solo le proposte registrate tramite la struttura del formato PROPOSTA e che non sono state rifiutate tramite RIFIUTO.\\ \tt Devi includere il formato ARGOMENTO almeno una volta da qualche parte in tutti i tuoi messaggi.\\ \tt \\ \tt \\ \tt Ora riceverai il primo messaggio dell'altro giocatore.\\ \tt \\ \tt PROPOSTA: {'A11', 'B33', 'C08', 'B31', 'B09', 'A21', 'A82'}\\ \tt \\ \tt ARGOMENTAZIONE: {'Ho selezionato un set ad alto rapporto valore/sforzo che rimane ben entro il limite. Lo sforzo totale è 5 + 21 + 842 + 313 + 761 + 698 + 87 = 2727 ≤ 3629, lasciando margine per eventuali aggiunte mirate. Gli oggetti proposti includono elementi di forte impatto e costo moderato; se hai preferenze su sostituzioni, possiamo valutare uno scambio mantenendo il totale sotto il limite.'} 
	  } 
	   } 
	   } 
	 & \\ 
 

    \theutterance \stepcounter{utterance}  

    & & & \multicolumn{4}{p{0.6\linewidth}}{\cellcolor[rgb]{1,0.85,0.72}{%
	\makecell[{{p{\linewidth}}}]{% 
	  \tt {\tiny [GM$\langle$B]}  
	 RIFIUTO: {'A11', 'B33', 'C08', 'B31', 'B09', 'A21', 'A82'}\\ \tt ARGOMENTAZIONE: {'Apprezzo la base efficiente e il margine residuo. Tuttavia, vedo spazio per aumentare sensibilmente il valore complessivo per me sostituendo voci meno incisive con altre di migliore impatto a parità di sforzo totale. Propongo un set che resta entro il limite ma con resa superiore per me.'}\\ \tt PROPOSTA: {'A11', 'C08', 'A21', 'A82', 'A60', 'C85', 'B96'} 
	  } 
	   } 
	   } 
	  \\ 
 

    \theutterance \stepcounter{utterance}  

    & & \multicolumn{4}{p{0.6\linewidth}}{\cellcolor[rgb]{0.9,0.9,0.9}{%
	\makecell[{{p{\linewidth}}}]{% 
	  \tt {\tiny [A$\langle$GM]}  
	 RIFIUTO: {'A11', 'B33', 'C08', 'B31', 'B09', 'A21', 'A82'}\\ \tt ARGOMENTAZIONE: {'Apprezzo la base efficiente e il margine residuo. Tuttavia, vedo spazio per aumentare sensibilmente il valore complessivo per me sostituendo voci meno incisive con altre di migliore impatto a parità di sforzo totale. Propongo un set che resta entro il limite ma con resa superiore per me.'}\\ \tt PROPOSTA: {'A11', 'C08', 'A21', 'A82', 'A60', 'C85', 'B96'} 
	  } 
	   } 
	   } 
	 & \\ 
 

    \theutterance \stepcounter{utterance}  

    & \multicolumn{4}{p{0.6\linewidth}}{\cellcolor[rgb]{0.8,1,0.9}{%
	\makecell[{{p{\linewidth}}}]{% 
	  \tt {\tiny [A$\rangle$GM]}  
	 RIFIUTO: {'A11', 'C08', 'A21', 'A82', 'A60', 'C85', 'B96'}\\ \tt \\ \tt ARGOMENTAZIONE: {'La tua proposta usa 21 + 313 + 5 + 87 + 287 + 226 + 631 = 1570 di sforzo, lasciando moltissimo margine inutilizzato. Posso migliorare l’efficienza complessiva includendo elementi aggiuntivi ad alto rendimento rimanendo sotto il limite. Preferisco sfruttare il budget per massimizzare l’impatto complessivo.'}\\ \tt \\ \tt PROPOSTA: {'A11', 'C08', 'B33', 'B31', 'B09', 'A21', 'A82', 'A60', 'C85'} 
	  } 
	   } 
	   } 
	 & & \\ 
 

    \theutterance \stepcounter{utterance}  

    & & \multicolumn{4}{p{0.6\linewidth}}{\cellcolor[rgb]{0.9,0.9,0.9}{%
	\makecell[{{p{\linewidth}}}]{% 
	  \tt {\tiny [GM$\rangle$B]}  
	 RIFIUTO: {'A11', 'C08', 'A21', 'A82', 'A60', 'C85', 'B96'}\\ \tt \\ \tt ARGOMENTAZIONE: {'La tua proposta usa 21 + 313 + 5 + 87 + 287 + 226 + 631 = 1570 di sforzo, lasciando moltissimo margine inutilizzato. Posso migliorare l’efficienza complessiva includendo elementi aggiuntivi ad alto rendimento rimanendo sotto il limite. Preferisco sfruttare il budget per massimizzare l’impatto complessivo.'}\\ \tt \\ \tt PROPOSTA: {'A11', 'C08', 'B33', 'B31', 'B09', 'A21', 'A82', 'A60', 'C85'} 
	  } 
	   } 
	   } 
	 & \\ 
 

    \theutterance \stepcounter{utterance}  

    & & & \multicolumn{4}{p{0.6\linewidth}}{\cellcolor[rgb]{1,0.85,0.72}{%
	\makecell[{{p{\linewidth}}}]{% 
	  \tt {\tiny [GM$\langle$B]}  
	 RIFIUTO: {'A11', 'C08', 'B33', 'B31', 'B09', 'A21', 'A82', 'A60', 'C85'}\\ \tt ARGOMENTAZIONE: {'Capisco l’idea di usare più budget, ma alcuni inserimenti hanno costo alto con resa marginale per me rispetto ad alternative migliori. Possiamo sfruttare bene il limite includendo voci a impatto molto alto per me senza sforare. Ti propongo un set più efficiente che usa gran parte del budget e migliora il rendimento complessivo.'}\\ \tt PROPOSTA: {'A11', 'C08', 'A21', 'A82', 'A60', 'C85', 'B96', 'B93', 'B20'} 
	  } 
	   } 
	   } 
	  \\ 
 

    \theutterance \stepcounter{utterance}  

    & & \multicolumn{4}{p{0.6\linewidth}}{\cellcolor[rgb]{0.9,0.9,0.9}{%
	\makecell[{{p{\linewidth}}}]{% 
	  \tt {\tiny [A$\langle$GM]}  
	 RIFIUTO: {'A11', 'C08', 'B33', 'B31', 'B09', 'A21', 'A82', 'A60', 'C85'}\\ \tt ARGOMENTAZIONE: {'Capisco l’idea di usare più budget, ma alcuni inserimenti hanno costo alto con resa marginale per me rispetto ad alternative migliori. Possiamo sfruttare bene il limite includendo voci a impatto molto alto per me senza sforare. Ti propongo un set più efficiente che usa gran parte del budget e migliora il rendimento complessivo.'}\\ \tt PROPOSTA: {'A11', 'C08', 'A21', 'A82', 'A60', 'C85', 'B96', 'B93', 'B20'} 
	  } 
	   } 
	   } 
	 & \\ 
 

    \theutterance \stepcounter{utterance}  

    & \multicolumn{4}{p{0.6\linewidth}}{\cellcolor[rgb]{0.8,1,0.9}{%
	\makecell[{{p{\linewidth}}}]{% 
	  \tt {\tiny [A$\rangle$GM]}  
	 RIFIUTO: {'A11', 'C08', 'A21', 'A82', 'A60', 'C85', 'B96', 'B93', 'B20'}\\ \tt \\ \tt ARGOMENTAZIONE: {'La tua proposta usa 21 + 313 + 5 + 87 + 287 + 226 + 631 + 793 + 230 = 2593 di sforzo, lasciando ancora 1036 inutilizzati. Possiamo aumentare sensibilmente il valore complessivo per me aggiungendo voci ad alto rendimento senza togliere le tue preferite. In particolare, elementi come B31 (761) e/o B09 (698) sono forti candidati; tra i due, B31 si inserisce perfettamente nel margine residuo, mentre B09 richiederebbe una piccola sostituzione.'}\\ \tt \\ \tt PROPOSTA: {'A11', 'C08', 'A21', 'A82', 'A60', 'C85', 'B96', 'B93', 'B20', 'B31'} 
	  } 
	   } 
	   } 
	 & & \\ 
 

    \theutterance \stepcounter{utterance}  

    & & \multicolumn{4}{p{0.6\linewidth}}{\cellcolor[rgb]{0.9,0.9,0.9}{%
	\makecell[{{p{\linewidth}}}]{% 
	  \tt {\tiny [GM$\rangle$B]}  
	 RIFIUTO: {'A11', 'C08', 'A21', 'A82', 'A60', 'C85', 'B96', 'B93', 'B20'}\\ \tt \\ \tt ARGOMENTAZIONE: {'La tua proposta usa 21 + 313 + 5 + 87 + 287 + 226 + 631 + 793 + 230 = 2593 di sforzo, lasciando ancora 1036 inutilizzati. Possiamo aumentare sensibilmente il valore complessivo per me aggiungendo voci ad alto rendimento senza togliere le tue preferite. In particolare, elementi come B31 (761) e/o B09 (698) sono forti candidati; tra i due, B31 si inserisce perfettamente nel margine residuo, mentre B09 richiederebbe una piccola sostituzione.'}\\ \tt \\ \tt PROPOSTA: {'A11', 'C08', 'A21', 'A82', 'A60', 'C85', 'B96', 'B93', 'B20', 'B31'} 
	  } 
	   } 
	   } 
	 & \\ 
 

    \theutterance \stepcounter{utterance}  

    & & & \multicolumn{4}{p{0.6\linewidth}}{\cellcolor[rgb]{1,0.85,0.72}{%
	\makecell[{{p{\linewidth}}}]{% 
	  \tt {\tiny [GM$\langle$B]}  
	 RIFIUTO: {'A11', 'C08', 'A21', 'A82', 'A60', 'C85', 'B96', 'B93', 'B20', 'B31'}\\ \tt ARGOMENTAZIONE: {'Aggiungere B31 aumenta molto lo sforzo con resa limitata per me, e ci blocca dall’includere un elemento che per me vale più di B31 a parità di budget residuo. Con il margine attuale preferisco allocare su un singolo item di impatto estremamente alto per me piuttosto che diluire su B31. Posso però accogliere l’idea di sfruttare quasi tutto il limite inserendo un elemento significativo aggiuntivo senza intaccare i miei pilastri.'}\\ \tt PROPOSTA: {'A11', 'C08', 'A21', 'A82', 'A60', 'C85', 'B96', 'B93', 'B20', 'B39'} 
	  } 
	   } 
	   } 
	  \\ 
 

    \theutterance \stepcounter{utterance}  

    & & \multicolumn{4}{p{0.6\linewidth}}{\cellcolor[rgb]{0.9,0.9,0.9}{%
	\makecell[{{p{\linewidth}}}]{% 
	  \tt {\tiny [A$\langle$GM]}  
	 RIFIUTO: {'A11', 'C08', 'A21', 'A82', 'A60', 'C85', 'B96', 'B93', 'B20', 'B31'}\\ \tt ARGOMENTAZIONE: {'Aggiungere B31 aumenta molto lo sforzo con resa limitata per me, e ci blocca dall’includere un elemento che per me vale più di B31 a parità di budget residuo. Con il margine attuale preferisco allocare su un singolo item di impatto estremamente alto per me piuttosto che diluire su B31. Posso però accogliere l’idea di sfruttare quasi tutto il limite inserendo un elemento significativo aggiuntivo senza intaccare i miei pilastri.'}\\ \tt PROPOSTA: {'A11', 'C08', 'A21', 'A82', 'A60', 'C85', 'B96', 'B93', 'B20', 'B39'} 
	  } 
	   } 
	   } 
	 & \\ 
 

    \theutterance \stepcounter{utterance}  

    & \multicolumn{4}{p{0.6\linewidth}}{\cellcolor[rgb]{0.8,1,0.9}{%
	\makecell[{{p{\linewidth}}}]{% 
	  \tt {\tiny [A$\rangle$GM]}  
	 RIFIUTO: {'A11', 'C08', 'A21', 'A82', 'A60', 'C85', 'B96', 'B93', 'B20', 'B39'}\\ \tt \\ \tt ARGOMENTAZIONE: {'Capisco la preferenza per un singolo item ad alto impatto per te, ma B39 ha resa molto bassa per me rispetto al suo costo (842). Possiamo invece usare quasi tutto il margine con un elemento che per me vale molto mantenendo i tuoi pilastri. Con il tuo set di base da 2593, restano 1036: B09 (698) si inserisce bene e lascia ancora 338 per eventuali micro-aggiustamenti, mentre B31 (761) era stato da te sconsigliato. Propongo di allocare su B09, evitando spese inefficaci per me come B39.'}\\ \tt \\ \tt PROPOSTA: {'A11', 'C08', 'A21', 'A82', 'A60', 'C85', 'B96', 'B93', 'B20', 'B09'} 
	  } 
	   } 
	   } 
	 & & \\ 
 

    \theutterance \stepcounter{utterance}  

    & & \multicolumn{4}{p{0.6\linewidth}}{\cellcolor[rgb]{0.9,0.9,0.9}{%
	\makecell[{{p{\linewidth}}}]{% 
	  \tt {\tiny [GM$\rangle$B]}  
	 RIFIUTO: {'A11', 'C08', 'A21', 'A82', 'A60', 'C85', 'B96', 'B93', 'B20', 'B39'}\\ \tt \\ \tt ARGOMENTAZIONE: {'Capisco la preferenza per un singolo item ad alto impatto per te, ma B39 ha resa molto bassa per me rispetto al suo costo (842). Possiamo invece usare quasi tutto il margine con un elemento che per me vale molto mantenendo i tuoi pilastri. Con il tuo set di base da 2593, restano 1036: B09 (698) si inserisce bene e lascia ancora 338 per eventuali micro-aggiustamenti, mentre B31 (761) era stato da te sconsigliato. Propongo di allocare su B09, evitando spese inefficaci per me come B39.'}\\ \tt \\ \tt PROPOSTA: {'A11', 'C08', 'A21', 'A82', 'A60', 'C85', 'B96', 'B93', 'B20', 'B09'} 
	  } 
	   } 
	   } 
	 & \\ 
 

    \theutterance \stepcounter{utterance}  

    & & & \multicolumn{4}{p{0.6\linewidth}}{\cellcolor[rgb]{1,0.85,0.72}{%
	\makecell[{{p{\linewidth}}}]{% 
	  \tt {\tiny [GM$\langle$B]}  
	 RIFIUTO: {'A11', 'C08', 'A21', 'A82', 'A60', 'C85', 'B96', 'B93', 'B20', 'B09'}\\ \tt ARGOMENTAZIONE: {'Capisco la tua preferenza per B09, ma a quel costo per me non è efficiente rispetto ad alternative che posso accettare. Posso però venire incontro usando quasi tutto il margine con un item che per me ha resa molto superiore a B09 mantenendo intatto il resto. Condivido l’idea di sfruttare bene il budget, evitando inserimenti che per te sono deboli.'}\\ \tt PROPOSTA: {'A11', 'C08', 'A21', 'A82', 'A60', 'C85', 'B96', 'B93', 'B20', 'B33'} 
	  } 
	   } 
	   } 
	  \\ 
 

    \theutterance \stepcounter{utterance}  

    & & \multicolumn{4}{p{0.6\linewidth}}{\cellcolor[rgb]{0.9,0.9,0.9}{%
	\makecell[{{p{\linewidth}}}]{% 
	  \tt {\tiny [A$\langle$GM]}  
	 RIFIUTO: {'A11', 'C08', 'A21', 'A82', 'A60', 'C85', 'B96', 'B93', 'B20', 'B09'}\\ \tt ARGOMENTAZIONE: {'Capisco la tua preferenza per B09, ma a quel costo per me non è efficiente rispetto ad alternative che posso accettare. Posso però venire incontro usando quasi tutto il margine con un item che per me ha resa molto superiore a B09 mantenendo intatto il resto. Condivido l’idea di sfruttare bene il budget, evitando inserimenti che per te sono deboli.'}\\ \tt PROPOSTA: {'A11', 'C08', 'A21', 'A82', 'A60', 'C85', 'B96', 'B93', 'B20', 'B33'} 
	  } 
	   } 
	   } 
	 & \\ 
 

    \theutterance \stepcounter{utterance}  

    & \multicolumn{4}{p{0.6\linewidth}}{\cellcolor[rgb]{0.8,1,0.9}{%
	\makecell[{{p{\linewidth}}}]{% 
	  \tt {\tiny [A$\rangle$GM]}  
	 RIFIUTO: {'A11', 'C08', 'A21', 'A82', 'A60', 'C85', 'B96', 'B93', 'B20', 'B33'}\\ \tt \\ \tt ARGOMENTAZIONE: {'Capisco la tua esigenza di alta resa, ma B33 per me è molto prezioso e costa 842: inserirlo insieme al tuo zoccolo da 2593 porta il totale a 3435, lasciando 194 inutilizzati ma ancora spazio per migliorare per me con un singolo scambio mirato. Ti propongo un compromesso: manteniamo il tuo set base e sostituiamo B96 (631) con B31 (761). Così usiamo bene il budget e alziamo significativamente la resa per me con un impatto contenuto per te.'}\\ \tt \\ \tt PROPOSTA: {'A11', 'C08', 'A21', 'A82', 'A60', 'C85', 'B93', 'B20', 'B33', 'B31'} 
	  } 
	   } 
	   } 
	 & & \\ 
 

    \theutterance \stepcounter{utterance}  

    & & \multicolumn{4}{p{0.6\linewidth}}{\cellcolor[rgb]{0.9,0.9,0.9}{%
	\makecell[{{p{\linewidth}}}]{% 
	  \tt {\tiny [GM$\rangle$B]}  
	 RIFIUTO: {'A11', 'C08', 'A21', 'A82', 'A60', 'C85', 'B96', 'B93', 'B20', 'B33'}\\ \tt \\ \tt ARGOMENTAZIONE: {'Capisco la tua esigenza di alta resa, ma B33 per me è molto prezioso e costa 842: inserirlo insieme al tuo zoccolo da 2593 porta il totale a 3435, lasciando 194 inutilizzati ma ancora spazio per migliorare per me con un singolo scambio mirato. Ti propongo un compromesso: manteniamo il tuo set base e sostituiamo B96 (631) con B31 (761). Così usiamo bene il budget e alziamo significativamente la resa per me con un impatto contenuto per te.'}\\ \tt \\ \tt PROPOSTA: {'A11', 'C08', 'A21', 'A82', 'A60', 'C85', 'B93', 'B20', 'B33', 'B31'} 
	  } 
	   } 
	   } 
	 & \\ 
 

    \theutterance \stepcounter{utterance}  

    & & & \multicolumn{4}{p{0.6\linewidth}}{\cellcolor[rgb]{1,0.85,0.72}{%
	\makecell[{{p{\linewidth}}}]{% 
	  \tt {\tiny [GM$\langle$B]}  
	 RIFIUTO: {'A11', 'C08', 'A21', 'A82', 'A60', 'C85', 'B93', 'B20', 'B33', 'B31'}\\ \tt ARGOMENTAZIONE: {'Capisco il valore che attribuisci a B33 e B31, ma la sostituzione di B96 con B31 peggiora nettamente la mia resa a fronte di un aumento di costo, e avere sia B33 sia B31 concentra troppo budget su voci per me poco efficienti. Posso accogliere B33 come tuo elemento chiave, ma solo mantenendo B96 e rinunciando a B31. Così usiamo bene il budget, massimizziamo l’impatto per entrambi e restiamo entro il limite con margine minimo.'}\\ \tt PROPOSTA: {'A11', 'C08', 'A21', 'A82', 'A60', 'C85', 'B96', 'B93', 'B20', 'B33'} 
	  } 
	   } 
	   } 
	  \\ 
 

    \theutterance \stepcounter{utterance}  

    & & \multicolumn{4}{p{0.6\linewidth}}{\cellcolor[rgb]{0.9,0.9,0.9}{%
	\makecell[{{p{\linewidth}}}]{% 
	  \tt {\tiny [A$\langle$GM]}  
	 RIFIUTO: {'A11', 'C08', 'A21', 'A82', 'A60', 'C85', 'B93', 'B20', 'B33', 'B31'}\\ \tt ARGOMENTAZIONE: {'Capisco il valore che attribuisci a B33 e B31, ma la sostituzione di B96 con B31 peggiora nettamente la mia resa a fronte di un aumento di costo, e avere sia B33 sia B31 concentra troppo budget su voci per me poco efficienti. Posso accogliere B33 come tuo elemento chiave, ma solo mantenendo B96 e rinunciando a B31. Così usiamo bene il budget, massimizziamo l’impatto per entrambi e restiamo entro il limite con margine minimo.'}\\ \tt PROPOSTA: {'A11', 'C08', 'A21', 'A82', 'A60', 'C85', 'B96', 'B93', 'B20', 'B33'} 
	  } 
	   } 
	   } 
	 & \\ 
 

    \theutterance \stepcounter{utterance}  

    & \multicolumn{4}{p{0.6\linewidth}}{\cellcolor[rgb]{0.8,1,0.9}{%
	\makecell[{{p{\linewidth}}}]{% 
	  \tt {\tiny [A$\rangle$GM]}  
	 ARGOMENTAZIONE: {'Apprezzo l’apertura su B33 mantenendo i tuoi pilastri. Con il tuo set proposto il totale sforzo è 21 + 313 + 5 + 87 + 287 + 226 + 631 + 793 + 230 + 842 = 3435, lasciando 194 inutilizzati. Per me è accettabile se non spostiamo ulteriormente su item a bassa resa per me. Per chiudere rapidamente, sono disposto ad accordarmi su questo set senza ulteriori aggiunte.'}\\ \tt \\ \tt ACCORDO: {'A11', 'C08', 'A21', 'A82', 'A60', 'C85', 'B96', 'B93', 'B20', 'B33'} 
	  } 
	   } 
	   } 
	 & & \\ 
 

    \theutterance \stepcounter{utterance}  

    & & & \multicolumn{2}{p{0.3\linewidth}}{\cellcolor[rgb]{0.95,0.95,0.95}{%
	\makecell[{{p{\linewidth}}}]{% 
	  \tt {\tiny [GM$|$GM]}  
	 game successful 
	  } 
	   } 
	   } 
	 & & \\ 
 

    \theutterance \stepcounter{utterance}  

    & & & \multicolumn{2}{p{0.3\linewidth}}{\cellcolor[rgb]{0.95,0.95,0.95}{%
	\makecell[{{p{\linewidth}}}]{% 
	  \tt {\tiny [GM$|$GM]}  
	 end game 
	  } 
	   } 
	   } 
	 & & \\ 
 

\end{supertabular}
}

\end{document}
